% BHU PYQ -- Transcribed & Corrected
% Paper: PHIMJ-41 / PHIMJ41: Contemporary Indian Philosophy
% Exam: B.A. (Hons.) Semester IV Examination, 2025-26 (NEP)
% Category: Major Core
% Department: Department of Philosophy, Faculty of Arts, Banaras Hindu University

\documentclass[11pt,a4paper]{article}
\usepackage[a4paper,top=2.5cm,bottom=2.5cm,left=2.8cm,right=2.8cm,headheight=14pt]{geometry}
\usepackage{amsmath,amssymb}
\usepackage[T1]{fontenc}
\usepackage[utf8]{inputenc}
\usepackage{lmodern,microtype,fancyhdr,enumitem,hyperref,lastpage,parskip}

\hypersetup{
    colorlinks=true,
    urlcolor=black,
    linkcolor=black,
    pdftitle={PHIMJ41: Contemporary Indian Philosophy -- BHU Sem IV 2025-26 (NEP)}
}

\pagestyle{fancy}
\fancyhf{}
\renewcommand{\headrulewidth}{0.4pt}
\renewcommand{\footrulewidth}{0.4pt}
\fancyhead[L]{\small \textit{Banaras Hindu University}}
\fancyhead[R]{\small \textit{PHIMJ-41: Contemporary Indian Philosophy (Sem IV 2025-26)}}
\fancyfoot[C]{\small Page \thepage\ of \pageref{LastPage}}

\newlist{parts}{enumerate}{2}
\setlist[parts,1]{label=(\alph*),leftmargin=*,itemsep=4pt,topsep=2pt}

\newcommand{\pts}[1]{\hfill \textit{[#1]}}

\begin{document}

\begin{center}
  {\large\bfseries BANARAS HINDU UNIVERSITY}\\[2pt]
  {\normalsize B.A. (Hons.) Semester IV Examination, 2025-26 (NEP)}\\[6pt]
  \rule{0.8\linewidth}{0.5pt}\\[6pt]
  {\large\bfseries DEPARTMENT OF PHILOSOPHY}\\[4pt]
  {\large\bfseries Paper Code: PHIMJ-41 : Contemporary Indian Philosophy}\\[4pt]
  {\normalsize Time: 2:30 Hours \hfill Max. Marks: 70}
\end{center}

\rule{\linewidth}{0.4pt}

\smallskip

\noindent \textit{(Write your Roll No. at the top immediately on the receipt of this question paper)}\\[2pt]
\noindent \textit{Note: This question paper comprises of three Sections 'A', 'B' and 'C'. Answer the questions as per the instructions given in each Section.}\\[1pt]
\noindent \textit{इस प्रश्न-पत्र में तीन खण्ड 'अ', 'ब' एवं 'स' हैं। प्रत्येक खण्ड में दिए गए निर्देशों के अनुसार प्रश्नों के उत्तर दीजिए।}

\bigskip

\begin{center}
  {\large\bfseries SECTION A / खण्ड-अ}\\[2pt]
  {\normalsize (Long Answer Type Questions / दीर्घ उत्तरीय प्रश्न)}
\end{center}

\noindent \textit{Note: Attempt any three questions out of the following five questions, each in 300 words. Each question carries 10 marks.}\\[2pt]
\noindent \textit{निम्नलिखित पाँच प्रश्नों में से किन्हीं तीन प्रश्नों का चयन कर, प्रत्येक का उत्तर 300 शब्दों में दीजिए। प्रत्येक प्रश्न 10 अंकों का है।}

\bigskip

\noindent \textbf{Question 1.} What are the main characteristics of contemporary Indian Philosophy? Explain it.\\[1pt]
समकालीन भारतीय दर्शन की मुख्य विशेषताये क्या हैं? व्याख्या कीजिए। \pts{10}

\medskip\rule{\linewidth}{0.2pt}\medskip

\noindent \textbf{Question 2.} How does Swami Vivekananda compare reality and God and give the different proofs for the existence of God? Elaborate it.\\[1pt]
स्वामी विवेकानन्द किस प्रकार सत् तथा ईश्वर की तुलना करते हैं तथा ईश्वर के अस्तित्व को सिद्ध करने हेतु विभिन्न तर्क देते हैं? सविस्तार बताइए। \pts{10}

\medskip\rule{\linewidth}{0.2pt}\medskip

\noindent \textbf{Question 3.} What are the different views expressed by M. K. Gandhi on Truth and how does he compare it with God? Describe it.\\[1pt]
सत्य पर एम० के० गाँधी द्वारा अभिव्यक्त विभिन्न विचार क्या हैं? तथा वे किस प्रकार इसकी तुलना ईश्वर से करते हैं? वर्णन कीजिए। \pts{10}

\medskip\rule{\linewidth}{0.2pt}\medskip

\noindent \textbf{Question 4.} Write an essay on Sri Aurobindo's Absolute.\\[1pt]
श्री अरविंद के निरपेक्ष पर एक निबंध लिखिए। \pts{10}

\medskip\rule{\linewidth}{0.2pt}\medskip

\noindent \textbf{Question 5.} What are the views of R. N. Tagore on his 'Doctrine of Maya'? Discuss it.\\[1pt]
आर० एन० टैगोर द्वारा 'माया' की अवधारणा पर व्यक्त किए गए विचार क्या हैं? विवेचना कीजिए। \pts{10}

\bigskip
\rule{\linewidth}{0.2pt}
\bigskip

\begin{center}
  {\large\bfseries SECTION B / खण्ड-ब}\\[2pt]
  {\normalsize (Short Answer Type Questions / लघु उत्तरीय प्रश्न)}
\end{center}

\noindent \textit{Note: Answer any four questions out of the following six questions, each in 150 words. Each question carries 5 marks.}\\[2pt]
\noindent \textit{निम्नलिखित छह प्रश्नों में से किन्हीं चार प्रश्नों का चयन कर, प्रत्येक का उत्तर 150 शब्दों में दीजिए। प्रत्येक प्रश्न 5 अंकों का है।}

\bigskip

\noindent \textbf{Question 6.} Clarify M. K. Gandhi's concept of means and end.\\[1pt]
एम० के० गाँधी के साधन तथा साध्य संबंधी अवधारणा को स्पष्ट कीजिए। \pts{5}

\medskip\rule{\linewidth}{0.2pt}\medskip

\noindent \textbf{Question 7.} Write a short note on Satyagrah.\\[1pt]
सत्याग्रह पर एक संक्षिप्त टिप्पणी लिखिए। \pts{5}

\medskip\rule{\linewidth}{0.2pt}\medskip

\noindent \textbf{Question 8.} What are the different aspects of Tagore's Humanism? Explain.\\[1pt]
टैगोर के मानवतावाद के विभिन्न पक्ष क्या हैं? व्याख्या कीजिए। \pts{5}

\medskip\rule{\linewidth}{0.2pt}\medskip

\noindent \textbf{Question 9.} How does S. Radhakrishnan describe the nature of intuitive apprehension? Describe it.\\[1pt]
एस० राधाकृष्णन 'अन्तर्दृष्टि' के स्वरूप का वर्णन किस प्रकार करते हैं? वर्णन कीजिए। \pts{5}

\medskip\rule{\linewidth}{0.2pt}\medskip

\noindent \textbf{Question 10.} What is the nature of self according to Mohammad Iqbal? Explain.\\[1pt]
मोहम्मद इकबाल के अनुसार 'आत्म' का स्वरूप क्या है? समझाइए। \pts{5}

\medskip\rule{\linewidth}{0.2pt}\medskip

\noindent \textbf{Question 11.} Write an essay on K. C. Bhattacharya's concept of Philosophy.\\[1pt]
के० सी० भट्टाचार्य के अनुसार दर्शन के स्वरूप पर एक निबंध लिखिए। \pts{5}

\bigskip
\rule{\linewidth}{0.2pt}
\bigskip

\begin{center}
  {\large\bfseries SECTION C / खण्ड-स}\\[2pt]
  {\normalsize (Very Short Answer Type Questions / अति लघु उत्तरीय प्रश्न)}
\end{center}

\noindent \textit{Note: Attempt all parts of this question, each maximum in two sentences. Each sub-question carries 2 marks.}\\[2pt]
\noindent \textit{इस प्रश्न के सभी भागों का उत्तर, अधिकतम दो वाक्यों में दीजिए। प्रत्येक उपप्रश्न 2 अंकों का है।}

\bigskip

\noindent \textbf{Question 12.}
\begin{parts}
  \item What is enjoying consciousness according to K. C. Bhattacharya?\\[1pt]
        के० सी० भट्टाचार्य के अनुसार उपभोगी चेतना क्या है? \pts{2}

  \item Which level of God has been explained as 'Supreme Truth Consciousness' by Sri Aurobindo?\\[1pt]
        श्री अरविंद सत्ता के किस स्तर को 'चरम सत्य चेतना' कहते हैं? \pts{2}

  \item Who is the author of 'Religion of Man'?\\[1pt]
        'रिलिजन ऑफ मैन' के रचयिता कौन हैं? \pts{2}

  \item Who has depicted Absolute as 'Pure Consciousness, Pure Freedom and Infinite Possibility'?\\[1pt]
        किसने निरपेक्ष सत् को 'शुद्ध चेतना', शुद्ध स्वतंत्रता तथा अनंत संभावना' कहा है? \pts{2}

  \item What Mohammad Iqbal means by 'Anal Haq'?\\[1pt]
        मोहम्मद इकबाल का 'अनल हक' से क्या तात्पर्य है? \pts{2}

  \item Mention the different types of Satyagrah.\\[1pt]
        सत्याग्रह के विभिन्न प्रकारों का उल्लेख कीजिए। \pts{2}

  \item Who has explained absolute as whole of perfection?\\[1pt]
        किसने निरपेक्ष को सर्वरूपेण पूर्ण कह कर व्याख्यायित किया है? \pts{2}

  \item What are the two possibilities that conscious force admits?\\[1pt]
        चित् शक्ति के स्वरूप में कौन-सी दो सम्भावनायें निहित हैं? \pts{2}

  \item Mention the names of eight different cords of being according to Sri Aurobindo.\\[1pt]
        श्री अरविंद के अनुसार सत्ता के विभिन्न 8 स्तरों का नामोल्लेख कीजिए। \pts{2}

  \item Which word has been used in Quran Sharif for 'Creativity'?\\[1pt]
        'सर्जनात्मकता' के लिए कुरान शरीफ में किस शब्द का प्रयोग किया गया है? \pts{2}
\end{parts}

\vfill
\rule{\linewidth}{0.4pt}

\begin{center}\small
\href{https://bhu-exam.vercel.app/}{Click here to visit our website}\\[4pt]
\href{https://me-aryan.vercel.app/}{Click here to visit our contributor}
\end{center}

\end{document}
